\documentclass[11pt]{article}
\usepackage{amsmath,amssymb,amsfonts}
\usepackage{geometry}
\usepackage{hyperref}

\geometry{margin=1in}

\begin{document}

\title{Color Image Processing Using Logarithmic Operations}
\author{Vasile P\^atr\c{a}scu, Vasile Buzuloiu}
\date{}
\maketitle

\begin{abstract}
In this paper, we present a mathematical model for color image processing. It is a logarithmic one. We consider the cube $(-1,1)^3$ as the set of values for the color space. We define two operations: addition $(+)$ and real scalar multiplication $(\times)$. With these operations the space of colors becomes a real vector space. Then, defining the scalar product $(\cdot)$ and the norm $\|\cdot\|$, we obtain a logarithmic Euclidean space. We show how we can use this model for color image enhancement and we present some experimental results.
\end{abstract}

\section{Introduction}

Image enhancement is an important branch of image processing. Its definition according to [1] can be expressed as follows: it consists of a collection of techniques that seek to improve the visual appearance of an image, or to convert the image to a form better suited for analysis by a human or a machine.

Another approach exists in the general setting of logarithmic representations suited for the imaging processes obtained by transmitted light or for human visual perception. Jourlin and Pinoli introduced a mathematical framework for this kind of nonlinear representation [2].

The logarithmic model for gray level images presented in [4,5] does work with bounded real sets. The extension of our model to color images is natural and we develop it here. We consider the unit cube of colors $(0,1)^3$ and transform it into the cube $(-1,1)^3$ by the simplest linear transformation, which translates the point $(0.5,0.5,0.5)$ to $(0,0,0)$. It is this new cube which plays the central role in our model.

\section{The Real Vector Space of the Colors}

We consider as the space of colors the cube
\[
E_3 = (-1,1)\times(-1,1)\times(-1,1).
\]
The three components of a vector $v\in E_3$ are noted $r$, $g$, and $b$ (red, green, blue). In the cube $E_3$ we define the addition $(+)$ and the real scalar multiplication $(\times)$.

\subsection{Addition}

For $v_1,v_2\in E_3$, with
\[
v_1=(r_1,g_1,b_1), \qquad v_2=(r_2,g_2,b_2),
\]
the sum $v_1(+ )v_2$ is defined by
\begin{equation}
v_1(+ )v_2=
\left(
\frac{r_1+r_2}{1+r_1r_2},
\frac{g_1+g_2}{1+g_1g_2},
\frac{b_1+b_2}{1+b_1b_2}
\right).
\end{equation}

The neutral element for addition is
\[
\theta=(0,0,0).
\]
Each element $v=(r,g,b)\in E_3$ has as its opposite the element
\[
-v=(-r,-g,-b).
\]
The addition is stable, associative, commutative, has a neutral element and each element has an opposite. It follows that this operation establishes on $E_3$ a commutative group structure.

The subtraction operation is defined by
\begin{equation}
v_1(-)v_2=
\left(
\frac{r_1-r_2}{1-r_1r_2},
\frac{g_1-g_2}{1-g_1g_2},
\frac{b_1-b_2}{1-b_1b_2}
\right).
\end{equation}

\subsection{Scalar Multiplication}

For $\lambda\in\mathbb{R}$ and $v=(r,g,b)\in E_3$, we define
\begin{equation}
\lambda(x)v=
\left(
\frac{(1+r)^\lambda-(1-r)^\lambda}{(1+r)^\lambda+(1-r)^\lambda},
\frac{(1+g)^\lambda-(1-g)^\lambda}{(1+g)^\lambda+(1-g)^\lambda},
\frac{(1+b)^\lambda-(1-b)^\lambda}{(1+b)^\lambda+(1-b)^\lambda}
\right).
\end{equation}

The two operations, addition $(+)$ and scalar multiplication $(\times)$, establish on $E_3$ a real vector space structure.

\subsection{The Euclidean Space of the Colors}

We define the scalar product $(\cdot)_{E_3}:E_3\times E_3\to\mathbb{R}$ by
\begin{equation}
(v_1\mid v_2)_{E_3}
= \varphi(r_1)\varphi(r_2)+\varphi(g_1)\varphi(g_2)+\varphi(b_1)\varphi(b_2),
\end{equation}
where
\[
\varphi:(-1,1)\to\mathbb{R}, \qquad \varphi(x)=\operatorname{arcth}(x).
\]

The norm is defined by
\begin{equation}
\|v\|_{E_3}=\sqrt{\varphi^2(r)+\varphi^2(g)+\varphi^2(b)}.
\end{equation}

Thus, our color space becomes a three-dimensional Euclidean space.

\section{The Vector Space of the Color Images}

A color image is a function defined on a bi-dimensional compact set $D\subset \mathbb{R}^2$ taking values in the color space $E_3$. We denote by $F(D,E_3)$ the set of color images defined on $D$.

\subsection{Addition}

For $f_1,f_2\in F(D,E_3)$ and $(x,y)\in D$,
\begin{equation}
(f_1(+ )f_2)(x,y)=f_1(x,y)(+ )f_2(x,y).
\end{equation}

The neutral element is the identically null function. The addition is stable, associative, commutative, has a neutral element and each element has an opposite. Therefore, this operation establishes on $F(D,E_3)$ a commutative group structure.

\subsection{Scalar Multiplication}

For $\lambda\in\mathbb{R}$, $f\in F(D,E_3)$ and $(x,y)\in D$,
\begin{equation}
(\lambda(x)f)(x,y)=\lambda(x)f(x,y).
\end{equation}

The two operations, addition $(+)$ and scalar multiplication $(\times)$, establish on $F(D,E_3)$ a real vector space structure.

\subsection{The Hilbert Space of the Color Images}

Let $f$ and $g$ be two integrable functions from $F(D,E_3)$. We define the scalar product by
\begin{equation}
(f_1\mid f_2)_{L^2(E_3)}=\int_D (f_1(x,y)\mid f_2(x,y))_{E_3}\,dx\,dy.
\end{equation}

The norm is then
\begin{equation}
\|f\|_{L^2(E_3)}=
\left(\int_D \|f(x,y)\|_{E_3}^2\,dx\,dy\right)^{1/2}.
\end{equation}

Thus, the color image space $F(D,E_3)$ becomes a Hilbert space.

\section{Applications in Enhancement of Color Images}

The enhancement algorithms presented in this paper consist of affine transformations
\begin{equation}
t:F(D,E_3)\to F(D,E_3), \qquad t(f)=\alpha(x)f(+)\beta(x)k,
\end{equation}
where $f$ is the original image and $h=t(f)$ is the enhanced image. Here $k$ is a constant image. The parameters $\alpha,\beta$ are determined so that the transformation satisfies some conditions.

Let $f:D\to E_3$ be a color image with scalar components
\[
f_r:D\to E,\qquad f_g:D\to E,\qquad f_b:D\to E.
\]

We consider the following subsets of the support $D$:
\[
D_{r1}=\{(x,y)\in D \mid f_r(x,y)<m(f_r(D))\},
\]
\[
D_{r2}=\{(x,y)\in D \mid f_r(x,y)>m(f_r(D))\},
\]
\[
D_{g1}=\{(x,y)\in D \mid f_g(x,y)<m(f_g(D))\},
\]
\[
D_{g2}=\{(x,y)\in D \mid f_g(x,y)>m(f_g(D))\},
\]
\[
D_{b1}=\{(x,y)\in D \mid f_b(x,y)<m(f_b(D))\},
\]
\[
D_{b2}=\{(x,y)\in D \mid f_b(x,y)>m(f_b(D))\},
\]
where
\begin{equation}
m(f_i(D))=\frac{1}{\operatorname{area}(D)}\int_D f_i(x,y)\,dx\,dy.
\end{equation}

In the discrete case, $\operatorname{area}(D)$ is replaced by $\operatorname{card}(D)$.

We define the following vectors:
\[
v_0=(m(f_r(D)),m(f_g(D)),m(f_b(D))),
\]
\[
v_1=(m(f_r(D_{r1})),m(f_g(D_{g1})),m(f_b(D_{b1}))),
\]
\[
v_2=(m(f_r(D_{r2})),m(f_g(D_{g2})),m(f_b(D_{b2}))),
\]
\[
u_L=(\operatorname{th}(1),\operatorname{th}(1),\operatorname{th}(1)),
\]
\[
w_0=(0,0,0), \qquad w_1=(-0.5,-0.5,-0.5), \qquad w_2=(0.5,0.5,0.5).
\]

\subsection{Algorithm A}

A constant translation on each component of the color. In this case we put
\begin{equation}
\forall x,y\in D,\qquad k(x,y)=u_L.
\end{equation}

We obtain the parameters $\alpha,\beta$ from the overdetermined linear system
\begin{equation}
\begin{cases}
\alpha(x)v_0(+)\beta(x)u_L=w_0,\\
\alpha(x)\big(v_1(-)v_0\big)=w_1,\\
\alpha(x)\big(v_2(-)v_0\big)=w_2.
\end{cases}
\end{equation}

Let us use the notations
\[
c_{vv}=\|v_0\|_{E_3}^2+\|v_1-v_0\|_{E_3}^2+\|v_2-v_0\|_{E_3}^2,
\]
\[
c_{uu}=\|u_L\|_{E_3}^2,\qquad
c_{vw}=(v_0\mid w_0)_{E_3}+(v_1\mid w_1)_{E_3}+(v_2\mid w_2)_{E_3},
\]
\[
c_{uv}=(v_0\mid u_L)_{E_3},\qquad
c_{uw}=(u_L\mid w_0)_{E_3},
\]
then the generalized solution, corresponding to the minimum mean square error, is
\begin{equation}
\alpha=\frac{c_{vw}c_{uu}-c_{uw}c_{uv}}{c_{vv}c_{uu}-(c_{uv})^2},
\end{equation}
\begin{equation}
\beta=\frac{-c_{vw}c_{uv}+c_{vv}c_{uw}}{c_{vv}c_{uu}-(c_{uv})^2}.
\end{equation}

The transformation is
\begin{equation}
t(f)=
\frac{c_{vw}c_{uu}-c_{uw}c_{uv}}{c_{vv}c_{uu}-(c_{uv})^2}\,x\,f
+
\frac{-c_{vw}c_{uv}+c_{vv}c_{uw}}{c_{vv}c_{uu}-(c_{uv})^2}\,x\,u_L.
\end{equation}

\subsection{Algorithm B}

A translation proportional to the mean. We suppose $v_0\in E_3$, $v_0\neq 0$. In this case we shall choose the constant vector
\begin{equation}
k=v_0.
\end{equation}

The overdetermined system from which the parameters $\alpha,\beta$ are found is
\begin{equation}
\begin{cases}
\alpha(x)v_0(+)\beta(x)v_0=w_0,\\
\alpha(x)v_1(+)\beta(x)v_0=w_1,\\
\alpha(x)v_2(+)\beta(x)v_0=w_2.
\end{cases}
\end{equation}

We again use the notations
\[
c_{vv}=\|v_0\|_{E_3}^2+\|v_1\|_{E_3}^2+\|v_2\|_{E_3}^2,
\]
\[
c_{uu}=3\|v_0\|_{E_3}^2,\qquad
c_{vu}=(v_0\mid v_0)_{E_3}+(v_1\mid v_0)_{E_3}+(v_2\mid v_0)_{E_3},
\]
\[
c_{vw}=(v_0\mid w_0)_{E_3}+(v_1\mid w_1)_{E_3}+(v_2\mid w_2)_{E_3},
\]
\[
c_{uw}=(v_0\mid w_0)_{E_3}+(v_0\mid w_1)_{E_3}+(v_0\mid w_2)_{E_3},
\]
and the minimum mean square error solution is
\begin{equation}
\alpha=\frac{c_{vw}c_{uu}-c_{uw}c_{vu}}{c_{vv}c_{uu}-(c_{vu})^2},
\end{equation}
\begin{equation}
\beta=\frac{-c_{vw}c_{vu}+c_{vv}c_{uw}}{c_{vv}c_{uu}-(c_{vu})^2}.
\end{equation}

The transformation is
\begin{equation}
t(f)=
\frac{c_{vw}c_{uu}-c_{uw}c_{vu}}{c_{vv}c_{uu}-(c_{vu})^2}\,x\,f
+
\frac{-c_{vw}c_{vu}+c_{vv}c_{uw}}{c_{vv}c_{uu}-(c_{vu})^2}\,x\,v_0.
\end{equation}

\section{Experimental Results}

The algorithm A gives
\[
\alpha=1.384,\qquad \beta=1.116,\qquad k=(0.762,0.762,0.762),
\]
\[
t(f)=1.384(x)f(+)\,1.116(x)k.
\]

The algorithm B gives
\[
\alpha=1.450,\qquad \beta=-1.447,\qquad k=(-0.694,-0.727,-0.580),
\]
\[
t(f)=1.450(x)f(-)1.447(x)k.
\]

For the boat image, the algorithm A gives
\[
\alpha=1.402,\qquad \beta=0.412,\qquad k=(0.762,0.762,0.762),
\]
\[
t(f)=1.402(x)f(+)\,0.412(x)k.
\]

The algorithm B gives
\[
\alpha=1.538,\qquad \beta=-1.941,\qquad k=(-0.194,-0.323,-0.338),
\]
\[
t(f)=1.538(x)f(-)1.941(x)k.
\]

\section{Conclusions}

Using this mathematical model we obtained two very simple algorithms for color image enhancement. The algorithm A improves the color contrast and saturation and preserves quite well the hue, while the algorithm B tends to attenuate the dominant colors which results in hue modifications. The main advantage consists in the ability to parameterize the algorithms and to simply introduce optimality criteria.

\begin{thebibliography}{9}
\bibitem{gonzales}
R.C. Gonzalez and P. Wintz, \emph{Digital Image Processing}, 2nd ed., Addison-Wesley, New York, 1987.

\bibitem{jourlin}
M. Jourlin and J.C. Pinoli, \emph{A model for logarithmic image processing}, Publication No. 3, Departement de Mathematiques, Univ. de Saint-Etienne, 1985.

\bibitem{oppenheim}
A.V. Oppenheim, \emph{Superposition in a class of non-linear systems}, Technical Report 432, Research Laboratory of Electronics, M.I.T. Cambridge, 1965.

\bibitem{patrascu2001}
V. P\^atr\c{a}scu and V. Buzuloiu, \emph{A mathematical model for logarithmic image processing}, SCI 2001, Orlando, USA, 2001.

\bibitem{patrascu2000}
V. P\^atr\c{a}scu and I. Voicu, \emph{An Algebraical Model for Gray Level Images}, OPTIM 2000, pp. 809--812, Bra\v{s}ov, Romania, 2000.

\bibitem{rosenfeld}
A. Rosenfeld and A.C. Kak, \emph{Digital Picture Processing}, Academic Press, New York, 1982.

\bibitem{stockham}
T.G. Stockham, \emph{Image processing in the context of visual model}, Proc. IEEE, Vol. 60, July 1972.
\end{thebibliography}

\end{document}